\documentclass[11pt]{article}
\usepackage[margin=1in]{geometry}
\usepackage[T1]{fontenc}
\usepackage[utf8]{inputenc}
\usepackage{lmodern}
\usepackage{microtype}
\usepackage{hyperref}
\usepackage{graphicx}
\usepackage{booktabs}
\usepackage{array}
\usepackage{enumitem}
\usepackage{xcolor}
\usepackage{amsmath}
\usepackage{amssymb}
\usepackage{float}

\hypersetup{
  colorlinks=true,
  linkcolor=blue,
  urlcolor=blue,
  citecolor=blue
}

\title{Language as a Protocol Layer for Specialist Financial Models:\\
LLM-Mediated DSL Routing for Regime-Aware Options and SPIKAN Backends}
\author{Draft for OSPIKAN Paper Development}
\date{\today}

\begin{document}
\maketitle

\begin{abstract}
Specialist quantitative models are often difficult to integrate into human-facing
decision workflows because they require custom interfaces, backend-specific data
plumbing, and repeated orchestration code. This problem is especially acute in
finance, where mathematically disciplined models must coexist with ambiguous
human requests, evolving market narratives, and auditable execution constraints.
We propose a language-mediated architecture in which a large language model
serves as a semantic coordination layer between human users and specialist
financial backends. Rather than allowing unconstrained natural-language control,
the system compiles user intent into a constrained domain-specific language
(DSL), which acts as the narrow waist between human reasoning and quantitative
execution.

We instantiate this architecture in a regime-aware options setting centered on
silver markets and a specialist backend called SPIKAN, a physics-informed
Kolmogorov-Arnold network designed for market shock and regime dynamics. The
repository already contains three major subsystems: a geometric market-state
pipeline for silver and macro data, an options DSL and KAN knowledge-store
stack, and a SPIKAN backend with PDE-informed training and outlook generation.
Our contribution is to reorganize these into a unified architecture with
explicit state contracts, option-data contracts, backend routing interfaces, and
reproducible evaluation paths.

The paper is designed to test three claims. First, the LLM can function as a
protocol compiler rather than a free-form assistant, translating natural
language into valid executable DSL with measurable reliability. Second, SPIKAN
can act as a specialist backend whose outputs remain grounded in financially
meaningful state variables, regime dynamics, and option-related decision tasks.
Third, a DSL-mediated router can reduce backend-integration complexity relative
to bespoke, backend-specific glue code. We define a benchmark program spanning
protocol quality, quantitative performance, calibration, stress-period
robustness, and integration cost. The resulting framework positions language not
as a replacement for quantitative modeling, but as an interoperability contract
between human intent and mathematically disciplined machine intelligence.
\end{abstract}

\section{Introduction}

Large language models are increasingly used as front ends for analytical and
decision-support systems. In many deployments, however, the LLM acts as an
unconstrained conversational wrapper around a collection of tools and models.
This pattern is flexible, but it is often brittle. The system may work for a
demo, yet remain difficult to validate, audit, extend, or compare across model
backends. These problems are amplified in financial applications, where user
requests are often ambiguous, model outputs must remain economically meaningful,
and the cost of silent failure is high.

This paper studies an alternative architecture. Instead of treating the LLM as
the decision engine, we treat it as a protocol layer between human language and
specialist quantitative models. The key mechanism is a constrained
domain-specific language that translates natural-language intent into typed,
executable commands. This design separates semantic coordination from numerical
inference. The LLM handles human-facing interpretation and query composition,
while the backend remains responsible for finance-native computation.

We instantiate this idea in a silver-market research platform that already
contains three substantial subsystems. The first is a multi-asset state
representation pipeline covering silver, macro, geometry, and regime features.
The second is an options DSL and KAN knowledge store for pricing, surfaces,
covariance, transitions, and what-if analysis. The third is SPIKAN, a
physics-informed Kolmogorov-Arnold network designed to model market shocks and
regime dynamics. The current repository demonstrates these components
individually, but the paper contribution is to reconstruct them around explicit
contracts and measurable evidence.

The business case is regime-aware options analysis. This domain is a useful test
bed because it forces interaction between ambiguous human questions, structured
contract parameters, regime-sensitive state variables, and specialist backend
reasoning. A human may ask for a stress-regime outlook, a repricing scenario, or
an options surface under changing macro conditions. A free-form LLM is
insufficient for such tasks because the outputs must remain executable and
auditable. A purely quantitative backend, on the other hand, is too rigid or
opaque for natural human interaction without substantial custom interface work.
The protocol-layer design aims to bridge that gap.

The paper advances three claims. First, the LLM can reliably compile natural
language into executable DSL under a controlled benchmark. Second, SPIKAN can
serve as a specialist backend whose outputs are quantitatively meaningful rather
than decorative. Third, the DSL-router architecture lowers integration cost when
comparing or onboarding backends. The emphasis is therefore not only predictive
performance, but also interoperability, auditability, and engineering economy.

\section{Architecture Overview}

The system is organized around six layers: a data plane, a state plane, a
protocol plane, a backend plane, an evaluation plane, and a product plane. The
data plane produces silver-centered market and macro features. The state plane
derives richer geometric and regime-sensitive representations, including Ricci
curvature, minimum-spanning-tree stress, geometric algebra rotor magnitude, and
manifold-tension diagnostics. The protocol plane converts natural language into
a constrained DSL. The backend plane executes requests through pricing, KAN, and
SPIKAN components. The evaluation plane measures protocol quality, backend
performance, and integration cost. The product plane exposes outputs through a
user-facing REPL and report artifacts.

\begin{figure}[H]
  \centering
  \fbox{\parbox{0.92\linewidth}{
    \centering
    Figure Placeholder: Human $\rightarrow$ LLM $\rightarrow$ DSL $\rightarrow$
    Router $\rightarrow$ Specialist Backend $\rightarrow$ Evidence-Aware Output
  }}
  \caption{System architecture showing the DSL as the narrow waist between human
  language and specialist quantitative backends. Replace with the final
  architecture figure during evidence binding.}
  \label{fig:architecture}
\end{figure}

\section{Canonical Contracts}

Two contracts anchor the design. The first is the market-state contract. The
current codebase uses state information in several different forms: DataFrame
columns in regime code, handcrafted dictionaries in the SPIKAN outlook adapter,
and substring-selected tensors in SPIKAN training. To remove that ambiguity, the
paper introduces a canonical \texttt{MarketStateSnapshot} with identity, base
market features, macro features, geometry features, regime features, fluid and
shock features, optional manifold-tension features, and provenance.

The second is the option-data contract. The options stack has a mature query
interface, but it lacks a canonical real-data option snapshot. The paper
therefore defines an \texttt{OptionSnapshot} with contract terms, underlying
reference fields, quote fields, derived fields such as moneyness and implied
volatility, and explicit quality flags. The joined benchmark unit is a
\texttt{StateOptionSnapshot}, formed by aligning option records with canonical
market state at the same evaluation time.

\begin{table}[H]
  \centering
  \caption{Contract summary for the narrow-waist data structures. Replace with
  the final contract table after implementation artifacts are fixed.}
  \label{tab:contracts}
  \begin{tabular}{p{0.22\linewidth} p{0.33\linewidth} p{0.33\linewidth}}
    \toprule
    Object & Core Fields & Primary Consumers \\
    \midrule
    MarketStateSnapshot & identity, macro, geometry, regime, tension, provenance & regime logic, SPIKAN, router, evaluation \\
    OptionSnapshot & contract terms, spot, quote, moneyness, implied vol, quality & pricing, KAN store, option benchmark \\
    BackendResult & query type, backend id, status, payload, provenance & router, formatter, benchmark harness \\
    \bottomrule
  \end{tabular}
\end{table}

\section{Protocol Layer}

The protocol layer is implemented as natural-language-to-DSL compilation with
validation-backed retries. Prompt rules constrain output to one command, enforce
capitalization and parameter formats, and provide few-shot examples across the
supported verb families. A validator checks whether the generated DSL is
syntactically and semantically executable. The paper reframes this layer as a
benchmarked protocol compiler rather than a convenience feature.

\begin{figure}[H]
  \centering
  \fbox{\parbox{0.92\linewidth}{
    \centering
    Figure Placeholder: natural-language request, generated DSL, normalized
    backend request, and structured result
  }}
  \caption{Protocol trace from human input to executable DSL and backend result.
  Replace with a real trace after protocol benchmark artifacts exist.}
  \label{fig:protocol-trace}
\end{figure}

\section{Specialist Backends}

The repository contains multiple backend families even before formal
modularization. A pricing path built around Black-Scholes with regime-conditioned
volatility from the KAN store supports option pricing, surfaces, and scenario
analysis. The KAN knowledge store also supports covariance and regime-transition
queries. SPIKAN consumes asset, macro, geometry, and time inputs to produce
dynamic signals and trader-facing outlook outputs. Existing regime and
Markov-style functions form natural baseline backends.

The paper defines a common backend interface over these families so that each
backend declares capabilities, executes a canonical request, and returns a
standardized result object with payload, status, provenance, and identity.

\section{Experimental Design}

The evaluation program contains three benchmark families: a protocol benchmark,
a quantitative benchmark, and an integration-cost study. The protocol benchmark
tests whether the LLM behaves like a reliable protocol compiler. The
quantitative benchmark tests whether SPIKAN provides meaningful value relative
to simpler baselines under equal data and time-split rules. The integration-cost
study tests whether the DSL-router architecture reduces backend onboarding and
workflow complexity relative to backend-specific glue.

\begin{table}[H]
  \centering
  \caption{Primary quantitative benchmark table placeholder. Replace with the
  first benchmark leaderboard once shared evaluation harnesses exist.}
  \label{tab:quant-benchmark}
  \begin{tabular}{lcccc}
    \toprule
    Backend & Task & Primary Metric & Calibration & Stress Slice \\
    \midrule
    SPIKAN & shock prediction & [TBD] & [TBD] & [TBD] \\
    Markov / HMM & shock prediction & [TBD] & [TBD] & [TBD] \\
    Tabular baseline & shock prediction & [TBD] & [TBD] & [TBD] \\
    Pricing baseline & option task & [TBD] & [TBD] & [TBD] \\
    \bottomrule
  \end{tabular}
\end{table}

\section{Expected Results Structure}

The final paper should report one table for protocol metrics, one leaderboard
table for quantitative results, one ablation table, one integration-cost table,
and a figure set covering architecture, protocol trace, state timeline,
quantitative comparison, and integration cost.

Current manuscript placeholders to replace during evidence binding:
\begin{itemize}[leftmargin=*]
  \item \texttt{[PROTOCOL\_VALID\_DSL\_RATE]}
  \item \texttt{[PROTOCOL\_SEMANTIC\_MATCH\_RATE]}
  \item \texttt{[SPIKAN\_PRIMARY\_TASK\_METRIC]}
  \item \texttt{[SPIKAN\_STRESS\_SLICE\_METRIC]}
  \item \texttt{[ABLATION\_NO\_GEOMETRY\_DELTA]}
  \item \texttt{[INTEGRATION\_COST\_ADAPTER\_LOC]}
  \item \texttt{[INTEGRATION\_COST\_FILES\_TOUCHED]}
\end{itemize}

\section{Discussion}

The main contribution of this work is architectural. Finance is a useful domain
because it makes the tradeoffs visible: free-form language is not safe enough,
but pure quantitative machinery is too rigid for broad interactive use without
substantial engineering. The DSL contract provides a narrow waist that absorbs
that tension. It gives the LLM a disciplined role and gives specialist models a
stable entry point into human-facing workflows.

The main limitation is that the strongest full-system claim still depends on
real-data option integration and benchmark completion. Until those are finished,
the paper should be described as a structured draft and evidence program rather
than a completed empirical result.

\section{Conclusion}

We have outlined a paper architecture in which language functions as a protocol
layer between humans and specialist financial models. In the proposed system,
the LLM does not replace quantitative reasoning. Instead, it compiles natural
language into a constrained DSL that can be routed through auditable, swappable,
finance-native backends. SPIKAN serves as the flagship specialist backend, while
the surrounding contracts and benchmarks make the broader architecture testable.

\appendix

\section{Minimal Figure Checklist}

The draft expects the following final figures:
\begin{enumerate}[leftmargin=*]
  \item architecture figure,
  \item state/regime/tension timeline,
  \item protocol trace figure,
  \item quantitative comparison figure,
  \item integration-cost figure.
\end{enumerate}

\end{document}
